\begin{textsample}{POS Dim 3 – mt – Score 29.00 – mt/mt\_inf\_24.txt}  \label{ex:f3_pos_075}
2.5.4 As Possibilidades de Experiências com \textbf{Bebês}, \textbf{Crianças} Bem Pequenas e \textbf{Crianças} Pequenas Pautados nos eixos interações e \textbf{brincadeiras} e nos direitos de aprendizagens e desenvolvimento, as experiências propostas às \textbf{crianças} da Educação \textbf{Infantil} Mato-Grossense, precisam garantir relações com os conhecimentos que compõem o patrimônio científico, ambiental e tecnológico, assim como, os saberes tradicionais da cultura local, e o respeito e \textbf{cuidado} com a sustentabilidade do planeta.
Assim, necessitam de oportunidades de : - Explorar e descobrir as propriedades dos objetos e materiais ( odor, cor, sabor, temperatura ) por meio de todos os sentidos ( olfato, paladar, audição, visão, tato ) ; - Experimentar as relações de causa e efeito ( transbordar, tingir, misturar, mover e remover etc. ) na interação com o mundo físico ; - \textbf{Manipular}, experimentar, arrumar e explorar o espaço por meio de experiências de deslocamentos de si e dos objetos ; - \textbf{Manipular} materiais diversos e variados para comparar as diferenças e semelhanças entre eles ; - Vivenciar diferentes ritmos, velocidades e fluxos nas interações e \textbf{brincadeiras} ( em danças, balanços, escorregadores etc. ) ; - Quantificar, \textbf{contar}, comparar, fazer cálculos, numerar, identificar numeração, fazer estimativas em relação à quantidade de pessoas ou objetos ; - Registrar quantidades, utilizando o traçado convencional ou não convencional, em situações significativas : pontuação de jogos, quantidades coletadas ou conquistadas ; - Comparar e classificar objetos com propriedades diversas : \textbf{peso} ( leve/pesado ), \textbf{volume} ( cheio/vazio ), espessura ( grosso/fino ), \textbf{textura} ( liso/áspero/macio ), cor e forma ; - Realizar atividades de \textbf{culinária} como receitas, envolvendo diferentes unidades de medidas : tempo de cozimento, quantidade de ingredientes, litro, quilograma, \textbf{colher}, xícara, entre outros ; - Amassar, transvazar, empilhar, encher, esvaziar, produzir \textbf{sons}, rolar objetos e materiais ; - Observar no meio natural e social as formas geométricas existentes, descobrindo semelhanças e diferenças entre objetos no espaço, combinando formas, estabelecendo relações espaciais e temporais, em situações que envolvam descrições orais, construções e representações ; - Explorar, orientar -se no espaço e indicar a posição de acordo com algumas relações : de vizinhança ( perto, longe, próximo ), de posição ( abaixo, acima, entre, ao lado, à direita, à esquerda ), de direção e sentido ( para a frente, para trás, para direita, para esquerda, para cima, para baixo, no mesmo sentido e em sentido diferente ) ; - Situar -se no espaço, indicando pontos de referência ; - Deslocar -se, em \textbf{brincadeiras} orientadas, verbalizando posições e distâncias nos percursos ; - Propiciar à \textbf{criança} o contato livre com diferentes materiais, portadores de atributos diversos, como cor, forma, tamanho, \textbf{textura}, temperatura, odor, utilidade entre outros, que possa estimular sua \textbf{percepção} e raciocínio ; - Participar de experiências em que o número tenha a função de memória de quantidade ; - Estimular a procura de objetos ou pessoas \textbf{escondidas} em diferentes lugares ; - Explorar espaços bidimensionais e tridimensionais utilizando materiais e ferramentas diferentes e construir conhecimentos sobre o \textbf{equilíbrio} das formas, \textbf{pesos} e tamanhos dos diferentes objetos que compõem seus primeiros jogos, com autonomia e independência ; - Realizar com a \textbf{criança} \textbf{pequenos} levantamentos de informações de assuntos de próprio interesse, organizando os dados em tabela ( por exemplo : \textbf{brinquedos} favoritos, cores prediletas, opções de lugares para fazer um \textbf{passeio}, opções de cardápio para um piquenique, etc. ) ; - Participar da elaboração de uma tabela para marcar as mudanças climáticas durante a semana ; - Representar objetos por meio de desenhos ou símbolos ; - Observar escrita numérica nos diferentes contextos em que se encontram ; - Utilizar circuitos numéricos para andar, pular, correr ; - Organizar espaços com \textbf{brinquedos} e objetos que contenha números como telefones, relógios, máquina de calcular e outros ; - Comparar quantidade com ou sem \textbf{ajuda} do professor ; - Identificar a passagem do tempo \textbf{apoiado} no \textbf{calendário} ; - Explorar com as \textbf{crianças} formas geométricas presentes nos ambientes, planificar e reconstruir \textbf{embalagens} ; - Incentivar a explicitação e representação da posição de pessoas e objetos, utilizando vocabulários pertinentes nos jogos, \textbf{brincadeiras} e nas situações em que essa ação seja necessária ; - Formar \textbf{pequenos} grupos de objetos pelas semelhanças e ainda, reunir elementos de duas ou mais coleções ; - Utilizar de noções simples de cálculo mental como estratégias para resolver problemas ; - Propor desafios para estimular a busca de estratégias próprias ; - Identificar os modos de ser, viver e trabalhar dos grupos sociais com os quais convive ; - Conhecer o próprio corpo, nomear algumas partes do mesmo e observar seu \textbf{crescimento} ; - Reconhecer diferenças e semelhanças entre sua organização familiar e a das outras \textbf{crianças} ; - Observar outras pessoas e comparar algumas de suas características pessoais respeitando as suas diferenças ; - Interagir com \textbf{adultos} e \textbf{crianças} de idades diferenciadas em \textbf{brincadeiras} que permitam o contato com a natureza e diferentes espaços. 2.5.5 Observação e Avaliação do Planejamento Para desenvolver experiências significativas e contextualizadas para as \textbf{crianças}, torna -se fundamental que o professor assuma postura de mediador.
Diante disso, ao planejar, é necessário pensar em importantes aspectos, tais como : - Encorajar as \textbf{crianças} a estabelecer diferentes relações com o mundo físico e social ; - Criar situações desafiadoras para as \textbf{crianças} pensarem sobre números, quantidades, formas, propriedade dos objetos, natureza e tecnologia ; - Propor questionamentos e dúvidas que mobilizem as \textbf{crianças} a indagarem sobre os aspectos do mundo físico, social e natural na construção de novos conhecimentos ; - Possibilitar às \textbf{crianças} levantar e interpretar suas hipóteses, considerando seus argumentos e ainda, oferecer meios para que elas ampliem suas informações e reformulem suas ideias iniciais, respeitando o patamar de sua lógica.
Lorenzato ( 2011 ), reitera que o professor possibilite muitas e distintas situações e experiências que devem pertencer ao mundo de vivência de quem vai construir sua própria aprendizagem.
Neste sentido, as \textbf{crianças} devem reproduzir ( escrevendo, falando, desenhando etc ) aquilo que aprenderam.
Algumas orientações O campo de experiências Espaços, tempos, quantidades, relações e transformações se faz presente na arte, na música, em histórias, na forma como se organiza o pensamento, nas \textbf{brincadeiras} e jogos \textbf{infantis}, na hora de dividir porções de lanche, de locomover no espaço, de organizar o tempo, ao comparar tamanho, distância, comprimento e formas, ou seja, as \textbf{crianças} da Educação \textbf{Infantil} nas situações \textbf{diárias} descobrem coisas iguais e diferentes, organizam, classificam e criam conjuntos, estabelecem relações, observam os tamanhos das coisas, \textbf{brincam} com as formas, ocupam um espaço e assim fazem diversificadas descobertas.
O educador ao propor suas ações didáticas, precisa se perguntar : - Como facilitar o desenvolvimento do pensamento \textbf{infantil}? - Quais são as primeiras explorações a serem realizadas?
Neste sentido as \textbf{crianças} da Educação \textbf{Infantil} Mato-Grossense, precisam ter asseguradas experiências que as levem a explorar e conhecer as diferentes formas e instrumentos do mundo físico e natural, assim como, apropriar -se das funções sociais desse campo, por meio dos jogos e \textbf{brincadeiras}, construindo formas convencionais e não convencionais de \textbf{registros}, cabendo ao professor fazer intervenções que envolvam a formulação de hipóteses, invenções, experimentos e pesquisas, dependendo da \textbf{curiosidade} e dos interesses das \textbf{crianças}.

% matched lemmas: adulto, ajuda, apoiar, bebê, brincadeira, brincar, brinquedo, calendário, colher, contar, crescimento, criança, cuidado, culinário, curiosidade, diário, embalagem, equilíbrio, esconder, infantil, manipular, passeio, pequeno, percepção, peso, registro, som, textura, volume
\end{textsample}
